\documentclass[
12pt,
openright,
twoside,
a4paper,
chapter=TITLE,
english,
portuguese % ou english
]{abntex2}

\usepackage[T1]{fontenc} 
\usepackage[utf8]{inputenc}
\usepackage[lmargin=3cm,tmargin=3cm,rmargin=2cm,bmargin=2cm]{geometry}
\usepackage[brazil]{babel}
\usepackage{times} %fonte Adobe Times Roman 
\usepackage{lastpage} % usado pela ficha catalográfica
\usepackage{color} % controle das cores		
\usepackage{graphicx} % inclusão de gráficos
\usepackage{float} % força as figuras a ficarem depois do texto que você colocou
\usepackage{chemfig,chemmacros}
\usepackage{upgreek}
\usepackage{microtype} % para melhorias de justificação	
\usepackage{pdfpages}
\usepackage{lastpage}
\usepackage{color}
\usepackage{makeidx}
\usepackage{quoting}
\usepackage{afterpage}
\usepackage{indentfirst} % indenta todos os primeiros parágrafos
\usepackage{setspace}
\usepackage{titlesec}
\renewcommand{\ABNTEXchapterfontsize}{\bfseries\normalsize}
\renewcommand{\ABNTEXsectionfontsize}{\bfseries\normalsize}
\renewcommand{\ABNTEXsubsectionfontsize}{\normalsize}
\usepackage{hyperref} % para permitir a criação de links
\usepackage[scaled]{helvet} % usa a fonte arial
\usepackage[T1]{fontenc} % seleção de fonte
\usepackage{fancyhdr}
\usepackage{verbatim}
\usepackage{amsmath,amssymb,amsfonts,textcomp} % vários pacotes juntos, utilizados para usar funções e outros símbolos matemáticos
\fancyhf{} % numeração no canto superior direito
\fancyheadoffset{0cm}
\renewcommand{\headrulewidth}{0pt} 
\renewcommand{\footrulewidth}{0pt}
\fancyhead[R]{\thepage}
\fancypagestyle{plain}{%
\fancyhf{}%
\fancyhead[R]{\thepage}%
}

\usepackage[alf]{abntex2cite} % Citações padrão ABNT

\renewcommand\familydefault{\sfdefault} % mantém a fonte como padrão do documento

% "Referências" em letras maiúsculas
\addto\captionsbrazil{%
  \renewcommand{\bibname}%
    {REFERÊNCIAS}%
}

\frenchspacing % Retira espaço obsoleto entre as frases.

\setlength\afterchapskip{18pt} % Ajusta o espaçamento entre o título do capitulo e o texto para 1,5

\begin{document}

% Capa
\input{capa}

% Folha de rosto: o aluno deve selecionar o arquivo da folha de rosto de acordo com seu tipo de trabalho. Basta trocar para uma das seguinte opções:
%folha de rosto_tcc
%folha de rosto_especializacao
%folha de rosto_dissertacao
%folha de rosto_tese
\input{folha de rosto_tcc}

% Ficha catalográfica
\includepdf{fichacatalografica_latex}

% Errata
\input{errata}
\cleardoublepage

% Folha de avaliação ou aprovação
\includepdf{folhadeavaliacao_latex}
\cleardoublepage

\input{dedicatoria}

% Agradecimentos
\input{agradecimentos}

% Epígrafe
\input{epigrafe}

% Resumo
\input{resumo}

% Abstract
\input{abstract}
\cleardoublepage

% Lista de ilustrações (elemento opcional - recomenda-se que sejam elaboradas a partir de 5 itens)
\listoffigures*
\cleardoublepage

% Lista de tabelas (elemento opcional - recomenda-se que sejam elaboradas a partir de 5 itens)
\listoftables*
\cleardoublepage

% Lista de abreviaturas e siglas (elemento opcional - recomenda-se que sejam elaboradas a partir de 5 itens)
\begin{siglas}
    \item[ABNT] Associa\c c\~ao Brasileira de Normas T\'ecnicas
    \item[abnTeX] ABsurdas Normas para TeX
	\item[EESC] Escola de Engenharia de S\~ao Carlos
	\item[IAU] Instituto de Arquitetura e Urbanismo
	\item[IBGE] Instituto Brasileiro de Geografia e Estat\'istica
	\item[ICMC] Instituto de Ci\^encias Matem\'aticas e de Computa\c c\~ao
	\item[IFSC] Instituto de F\'isica de S\~ao Carlos
	\item[IQSC] Instituto de Qu\'imica de S\~ao Carlos
	\item[PDF] Portable Document Format
	\item[TCC] Trabalho de Conclus\~ao de Curso
	\item[USP] Universidade de S\~ao Paulo
	\item[USPSC] Campus USP de S\~ao Carlos    
\end{siglas}

% Lista de símbolos (elemento opcional - recomenda-se que sejam elaboradas a partir de 5 itens)
\begin{simbolos}
  \item[$ \Gamma $] Letra grega Gama
  \item[$ \Lambda $] Lambda
  \item[$ \zeta $] Letra grega min\'uscula zeta
  \item[$ \in $] Pertence
\end{simbolos}

% Sumário
\pdfbookmark[0]{\contentsname}{toc}
\tableofcontents*
\cleardoublepage

\textual
\pagestyle{simple} % tira cabeçalho com nome do capítulo e exibe só paginação

% Capítulo 1
\include{cap1-introducao}

% Capítulo 2
\input{cap2-revisao bibliografica}

% Capítulo 3
\input{cap3-modelos de citacao}

% Capítulo 4
\input{cap4-conclusao}

\bibliography{referencias}

% Glossário
\input{glossario}

% Apêndice
\input{apendice}

% Anexo
\input{anexo}

% Índice
\input{indice}

\end{document}
